\documentclass[11pt,a4paper]{article}
% Shared preamble for RuPhO English problem statements (match T1 2026 style).
% Use: \input{latex/rupho_en_preamble.tex} after \documentclass.
\usepackage[margin=1in]{geometry}
\usepackage{amsmath,amssymb}
\usepackage{graphicx}
\usepackage{float}
\usepackage{enumitem}
\usepackage{microtype}
\usepackage{caption}
\usepackage{tabularx}
\usepackage{hyperref}

\captionsetup{font=small,labelfont=bf,skip=6pt}

\setlist[itemize]{leftmargin=1.2cm}
\setlist[enumerate]{leftmargin=1.2cm}

% One outer box per subproblem: [id | statement | points] (no inner rules)
\newcommand{\problembox}[3]{%
  \par\addvspace{0.42em}%
  \noindent
  \setlength{\fboxsep}{7.5pt}%
  \setlength{\fboxrule}{0.95pt}%
  \fbox{%
    \begin{minipage}{\dimexpr\linewidth-2\fboxsep-2\fboxrule\relax}
    \begin{tabularx}{\linewidth}{@{}>{\raggedright\arraybackslash\bfseries}p{2.1em} @{\hspace{0.9em}} >{\raggedright\arraybackslash}X @{\hspace{0.9em}} >{\raggedleft\arraybackslash\bfseries}p{2.4em}@{}}
    #1 & #2 & #3
    \end{tabularx}
    \end{minipage}%
  }%
  \par\addvspace{0.32em}%
}


\title{\textbf{Detail}}
\author{\normalsize Y14 (2014) --- English translation}
\date{}
\begin{document}
\maketitle

A part with mass $m$ can move without friction along vertical guides and rotate freely around a horizontal axis perpendicular to the guides and the drawing plane. The part has a through hole into which a smooth light rod is inserted, one end of which is hinged on a horizontal plane at a distance $L$ from the base of the vertical guides (Fig. 1). Between the part and the hinge, a spring attached to them with a stiffness coefficient $k$ is put on the rod. Free spring length $3L$. A light thread is tied to the part, which in the initial state is tensioned by a vertically upward force $F=mg$.

A part with mass $m$ can move without friction along vertical guides and rotate freely around a horizontal axis perpendicular to the guides and the drawing plane. The part has a through hole into which a smooth light rod is inserted, one end of which is hinged on a horizontal plane at a distance $L$ from the base of the vertical guides (Fig. 1). Between the part and the hinge, a spring attached to them with a stiffness coefficient $k$ is put on the rod. Free spring length $3L$. A light thread is tied to the part, which in the initial state is tensioned by a force $F=mg$ directed vertically upward.

\begin{figure}[H]
    \centering
    \includegraphics[width=0.72\textwidth]{T1_EN_assets/fig1.jpg}
    \caption{Figure 1.}
\end{figure}

\problembox{A1}{What should be the minimum mass $m_\text{min}$ of the part so that with a slow (quasi-static) decrease in force $F$, the part can move the rod to a horizontal position under the influence of gravity. Ignore friction.}{0}

If the force $F$ is reduced to zero instantly, then the part can move the rod to a horizontal position with a smaller mass $m_1$. The graph shows the dependence of the projection of the elastic force on the vertical axis, in relative units, on the distance from the load to the horizontal plane.

\begin{figure}[H]
    \centering
    \includegraphics[width=0.72\textwidth]{T1_EN_assets/fig2.jpg}
    \caption{Figure 2.}
\end{figure}

\problembox{A2}{Determine the mass of $m_1$ using the graph provided.}{0}

\problembox{A3}{The force $F$ was eliminated (the thread was cut), and the mass load $m_1$ was transferred to a stable equilibrium position. Determine the period of small oscillations near this position.}{0}

\vspace{1em}
\noindent\textit{Source:} \url{https://pho.rs/p/3963}

\end{document}